\label{sec:literature}

Our analysis connects sports officiating research, home-advantage and crowd-pressure studies, play-by-play sports analytics, and the sequential-decision / gambler's-fallacy literature.

\subsection{Officiating bias and home advantage}

\citet{pricewolfers2010} show that NBA foul calls vary with referee--player racial composition, establishing that whistle outcomes are measurable and policy-relevant. Related NBA work uses play-by-play data to study home-court favoritism, game-state incentives, and foul outcomes \citep{price2012}. More broadly, studies document home advantage in penalties and scoring, social pressure on officials \citep{garicano2005,dawson2010,buraimo2010}, and crowd-related distortions in refereeing decisions \citep{page2010}. Popular and academic accounts also emphasize home-court and officiating narratives in professional basketball \citep{moskowitz2011}. Prior work studies referee bias and officiating outcomes, but less attention has been paid to within-game sequential dependence in foul-call \emph{direction}. We model whether the \emph{next} foul falls on home or away conditional on explicit pre-call state. Unconditional foul rates in our main sample are near even ($P(\text{foul against home})\approx 0.50$), so our contribution is conditional sequential dependence, not a constant home tilt.

\subsection{Sequential decisions and the law of small numbers}

Psychologists document over-inference from short sequences \citep{tversky1971}. \citet{pope2011} and \citet{chen2016} find negative serial correlation in high-stakes decisions by athletes and professional decision-makers. Sports-officiating studies provide closer parallels: experimental work on soccer penalty decisions documents negative serial correlation in calls against the same team \citep{plessner2001}; college basketball officials tend to call fouls on the team with fewer fouls, evening out foul differentials \citep{anderson2009}; and NBA referees show increased scrutiny on one team after a potentially difficult call on the other \citep{gift2015}. We adopt neutral language---\emph{sequential dependence}---and implement a within-game shuffle placebo designed to falsify order-specific associations while preserving foul totals. This falsification logic is central to our design: not all regressors should respond to order destruction.

\subsection{Play-by-play sports econometrics}

Granular NBA play-by-play enables conditioning on score, clock, possession, and period. Foundational work in basketball analytics formalizes box-score and play-by-play metrics \citep{kubatko2007}; subsequent models use event-level and player-tracking data to predict possession outcomes and value player actions \citep{cervone2014ssac,cervone2016}. Prior PBP work emphasizes shot efficiency, lineup value, and strategic fouling; fewer papers treat foul events as outcomes with sequential state. We contribute a transparent pipeline from raw PBP to analysis-ready foul events, with validation against box-score foul totals (Appendix~\ref{app:validation}).

\subsection{Gap and positioning}

Few studies jointly model previous call direction, foul-count state, game fixed effects, and within-game placebo falsification on a multi-season NBA foul-event panel. Existing literature studies referee bias, home advantage, foul-count leveling, and event-level basketball analytics; this paper fills the gap by testing order-specific sequential dependence in foul-call direction with an explicit falsification design. We further show that offensive fouls---excluded from the main sample---obey different assignment rules, defining an important boundary for pooled interpretations.
